% ===========================================================================
% Poste 15 — Crédit pour la TPS/TVH
% ===========================================================================
\poste{15}{Crédit pour la TPS/TVH}{CA\_tps}

\pdesc Crédit fédéral remboursable et non imposable, versé trimestriellement
par l'Agence du revenu du Canada pour compenser la taxe sur les produits et
services payée par les ménages à revenu faible ou modeste.

\pobj Atténuer la régressivité de la taxe de vente : rembourser aux ménages
modestes une partie de la TPS qu'ils paient sur leur consommation.

\pregle Le crédit additionne un montant de base par adulte et par enfant et,
selon la composition, l'un de deux suppléments, puis applique une réduction
de \pc{5} sur le revenu familial net rajusté (équation~\ref{eq:afni})
au-delà du seuil $s_r$ :
\begin{align*}
\text{base} &= n_{ad} \times B + n \times E \\
\text{supplément} &= \begin{cases}
S_{mono} & \text{famille monoparentale (équivalent de conjoint)} \\
\min\bigl(0{,}02 \times \pos{\AFNI - s_p},\; S_{max}\bigr) &
\text{personne seule sans enfant} \\
0 & \text{couple}
\end{cases} \\
\text{crédit} &= \pos{\text{base} + \text{supplément}
 - 0{,}05 \times \pos{\AFNI - s_r}}
\end{align*}
Particularité du supplément \og personne seule \fg{} : il \emph{croît} avec
le revenu (\pc{2} du revenu au-delà de $s_p$, jusqu'à $S_{max}$) avant que
la réduction générale ne s'applique --- un profil en tente, en miniature.
L'année affichée correspond à l'année de versement (juillet à juin).

\pparams

\begin{center}
\small
\begin{tabular}{l r r}
\toprule
Paramètre & 2025 & 2026 \\
\midrule
Montant de base par adulte ($B$)            & \D{349}   & \D{356} \\
Montant par enfant ($E$)                    & \D{184}   & \D{187} \\
Supplément monoparental ($S_{mono}$)        & \D{349}   & \D{356} \\
Supplément personne seule --- seuil ($s_p$) & \D{11337} & \D{11564} \\
Supplément personne seule --- taux          & \pc{2}    & \pc{2} \\
Supplément personne seule --- maximum ($S_{max}$) & \D{184} & \D{187} \\
Seuil de réduction ($s_r$)                  & \D{45521} & \D{46432} \\
Taux de réduction                           & \pc{5}    & \pc{5} \\
\bottomrule
\end{tabular}
\end{center}

\pnotes Trois précisions.

(1) \emph{Allocation canadienne pour l'épicerie et les besoins essentiels
--- divergence connue.} À compter de juillet 2026, le crédit pour la
TPS/TVH est remplacé par l'\emph{Allocation canadienne pour l'épicerie et
les besoins essentiels}, dont les montants sont bonifiés de \pc{25} pour
cinq ans (Lois annuelles du Canada 2026, ch.~1)\footnote{Texte de loi :
\url{https://laws-lois.justice.gc.ca/fra/LoisAnnuelles/2026_1/page-1.html}.
Un versement unique de rattrapage (environ \pc{50} de la valeur 2025-2026),
payé en juin 2026, est un paiement ponctuel hors période de versement, non
modélisé.}. La colonne 2026 du tableau (période juillet 2026--juin 2027)
reproduit le calculateur officiel du MFQ dans son édition de décembre 2025,
\emph{antérieure à cette loi} : les montants n'incluent pas la bonification
(par exemple \D{356} par adulte plutôt que \D{445}). Conformément à la
règle de parité du projet (section~\ref{sec:validation}), cette divergence
est documentée plutôt que corrigée ; elle sera résorbée à la prochaine
édition du calculateur officiel.

(2) \emph{Paramétrisation.} L'article~122.5 LIR s'énonce autrement que le
modèle : le premier enfant d'un parent seul y reçoit le montant
\og équivalent de conjoint \fg{} ($B$ au lieu de $E$) et le supplément de
célibataire est accordé \emph{en entier} au parent seul, sans accumulation.
La forme du modèle --- $n \times E$ plus un supplément monoparental égal à
$B$ --- est algébriquement équivalente, car le supplément maximal vaut
exactement $E$ : $B + B + (n-1) \times E + S_{max} = B + n \times E + B$.
Les deux écritures donnent le même crédit pour toute famille monoparentale.

(3) La \emph{garde partagée} (la moitié du crédit de l'enfant à chaque
parent) n'est pas couverte par les entrées du modèle.

% ============================================================================
% GÉNÉRÉ par scripts/exemples-guide.ts — NE PAS ÉDITER À LA MAIN.
% Régénérer : npm run exemples (chaque chiffre est vérifié contre le moteur).
% ============================================================================
\pexemples Montants de base par adulte (\D{349}) et par enfant
(\D{184}), suppléments selon la composition, réduction de
\pc{5} de l'$\AFNI$ au-delà de \D{45521}.

\medskip\noindent\textbf{M3 --- personne vivant seule, 30~ans, revenu de travail de \D{15000}.}
Personne seule sans enfant : le \emph{supplément pour célibataire}
s'accumule à \pc{2} de l'$\AFNI$ au-delà de \num{11337},
plafonné à \num{184}.
\begin{align*}
\text{supplément} &= \min\bigl(\pc{2} \times (\num{14885} - \num{11337}),\ \num{184}\bigr) = \num{70.96} \\
\text{crédit} &= \num{349} + \num{70.96} - 0 = \num{419.96}
\end{align*}
Crédit TPS : \D{419.96}.

\medskip\noindent\textbf{M6 --- famille monoparentale, 35~ans, revenu de travail de \D{35000}, un enfant : 3~ans (garde subventionnée, \D{2000}).}
Famille monoparentale : base adulte \num{349} $+$ enfant
\num{184} $+$ supplément monoparental \num{349} $=$
\num{882}. L'$\AFNI$ (\num{32685}) est sous le seuil de
réduction : crédit \emph{plein} de \D{882}.

\medskip\noindent\textbf{M4 --- personne vivant seule, 40~ans, revenu de travail de \D{50000}.}
\begin{align*}
\text{total avant réduction} &= \num{349} + \num{184} = \num{533} \\
\text{réduction} &= \pos{\num{49535} - \num{45521}} \times \pc{5} = \num{200.70} \\
\text{crédit} &= \pos{\num{533} - \num{200.70}} = \num{332.30}
\end{align*}
Le supplément pour célibataire est au plafond (\num{184}), mais la
réduction mord : crédit de \D{332.30}.

\medskip\noindent\textbf{M5 --- personne vivant seule, 45~ans, revenu de travail de \D{100000}.}
La réduction (\num{2670.25}) dépasse le total : crédit
\emph{éteint} (\D{0}).


\prefs Loi de l'impôt sur le revenu (L.R.C. 1985, ch.~1
(5\ieme{}~suppl.)), art.~122.5 ; Agence du revenu du Canada ; CFFP,
\emph{Guide des mesures fiscales}, fiche \og Crédit d'impôt pour la
TPS/TVH \fg{}\footnote{\url{https://cffp.recherche.usherbrooke.ca/outils-ressources/guide-mesures-fiscales/credit-impot-tps-tvh/}}.
